\documentclass[11pt]{article}
\usepackage[T1]{fontenc}
\usepackage[margin=1in]{geometry}
\usepackage{amsmath,amssymb,graphicx,booktabs,longtable,array}
\usepackage[hidelinks]{hyperref}
\newcommand{\pending}[1]{\par\noindent\textbf{UNFILLED:} #1\par}
\title{Phaseward: a research-paper scaffold for coherent one-dimensional transport}
\author{Author names and affiliations: to be supplied}
\date{Draft scaffold; not a report of a discovery}
\begin{document}
\maketitle
\begin{abstract}
\pending{Write the abstract only after the frozen study, numerical validation,
full literature review and limitations have been assessed. No performance,
novelty, hardware validation or discovery claim is asserted here.}
\end{abstract}
\section{Introduction}
\pending{State the transport question, motivation and boundaries of the idealized
model. Make any proposed contribution conditional on evidence. Distinguish
simulation from experiment and gate energies from physical electrode voltages.}
\section{Related Work}
The screened bibliography includes scattering/conductance theory
\cite{fisher1981}, disorder-aware transport control \cite{coopmans2022},
derivative-based robust control \cite{propson2022}, disorder-averaged control
\cite{araki2023}, and perturbative robustness \cite{poggi2024}.
These are starting points for comparison, not a claim that their models,
observables or constraints coincide with Phaseward.
\pending{Complete full-text review using the literature matrix. Establish any
precise gap with evidence; absence from this initial screening proves no gap.}
\section{Model}
The assumed physical model is a coherent single electron with constant effective
mass $m^*>0$ in one spatial dimension:
\[
i\hbar\partial_t\psi(x,t)=-\frac{\hbar^2}{2m^*}\partial_x^2\psi(x,t)+V(x)\psi(x,t).
\]
Here $x$ is position (internally metres), $t$ is time (seconds), $\hbar$ is the
reduced Planck constant (joule seconds), $V$ is potential energy (joules), and
$\psi$ is a position amplitude (metres$^{-1/2}$) for a normalized packet.
The domain has a finite device region and uniform semi-infinite leads. The
public API uses nm, meV and ps with explicit conversion. This equation is an
assumption, not a derivation from classical mechanics.

For elastic propagation, let $t(E)$ and $r(E)$ be complex amplitudes in the
unit-flux incoming/outgoing lead basis at explicitly fixed reference planes.
Then $T(E)=|t(E)|^2$ and $R(E)=|r(E)|^2$ are dimensionless probabilities for the
open channel. Phase is retained in the amplitudes. Let $A(E)$ be an
energy-normalized incoming amplitude (meV$^{-1/2}$ at the API), with
$\int_{\mathcal E}|A(E)|^2\,dE=1$ over the declared finite interval
$\mathcal E$ of propagating lead energies. Define
\[
P=\int_{\mathcal E}|A(E)t(E)|^2\,dE,\qquad
b(E)=\frac{A(E)t(E)}{\sqrt P},\qquad
F=\left|\int_{\mathcal E}b_{\rm target}(E)^*b(E)\,dE\right|^2.
\]
$b$ and $b_{\rm target}$ are separately normalized conditional transmitted
states; $F$ is undefined when either cannot be normalized. $P$ and $F$ must be
reported separately. A global phase cancels from $F$. Optional delay alignment
uses a declared bounded lag search; its bound, sign, resolution and boundary
status accompany raw $F$. Numerical quadrature and finite-window normalization
are approximations with recorded settings.

\section{Numerical Methods}
\pending{Describe the centered finite-difference Hamiltonian and exact retarded
semi-infinite discrete leads, open-band restrictions, pivoted banded solves,
flux normalization and phase reference planes. Report spatial and energy
convergence, residuals, analytical checks and the second-method cross-check.
Copy actual settings from the support tables. Never equate machine precision
conservation with continuum accuracy.}
\section{Optimization}
For a channel coordinate $x$ in nm, write the potential energy in meV as
\[
V(x;u,\eta)=V_0(x)+\sum_{j=1}^M u_jg_j(x)+\sum_{k=1}^K\eta_k\phi_k(x).
\]
$V_0$ is the declared background, $u_j$ are bounded gate energy amplitudes (meV),
and $g_j$ are dimensionless peak-normalized shapes confined to the channel.
The covariance shapes obey $\Phi^\mathsf{T}W\Phi=I$, where $W$ contains spatial
trapezoidal weights (nm); hence $\phi_k$ has units nm$^{-1/2}$.
For independent standard-normal $z_k$, $\eta_k=\sqrt{\lambda_k}z_k$ has
units meV$\sqrt{\text{nm}}$, and its variance $\lambda_k$ has units meV$^2$nm.
$K$ is the retained covariance rank. This weighted finite-grid representation
approximates the declared spatial covariance; its omitted variance is reported.

The candidate surrogate is
\[
J(u)=\sum_{k=1}^K\lambda_k
\int_{\mathcal E}w(E)
\left|\left.\frac{\partial t(E;u,\eta)}{\partial\eta_k}\right|_{\eta=0}\right|^2dE.
\]
$w\geq0$ is a normalized importance density (meV$^{-1}$), normally $|A|^2$,
with any alternative declared. $J$ is dimensionless; $J/\sigma^2$ has units
meV$^{-2}$ for nonzero nominal disorder amplitude $\sigma$ in meV. The derivative
acts on the complex amplitude at fixed leads and phase planes. To first order
in zero-mean covariance coefficients, $J$ describes weighted mean squared
amplitude perturbation. It is neither conditional fidelity nor a failure
probability, and ignores higher-order response and omitted covariance modes.
\pending{Insert actual differentiation steps, step checks, optimizer settings,
constraints and costs from the frozen run. Compare fixed gates, clean optimization,
training-set Monte Carlo optimization and this surrogate with fair budgets.}

\section{Experimental Protocol}
\pending{Declare the frozen target, clean-throughput and raw-fidelity constraints,
gate bounds, A/B/C/D baselines, common starts and work caps, seed partitions,
primary estimand, practical margin, multiplicity, and failed-trial policy.
Separate numerical diagnostics from physical findings. An existing observed run
cannot be retrospectively preregistered. Include training/initialization
replications if the claimed estimand extends beyond the specific frozen designs.}
\section{Results}
\pending{Results deliberately unfilled. Do not invent measurements, improvement
percentages, significance, discoveries or favorable examples. A completed run
alone is not publication readiness. After experiment and numerical review,
authors may write a Results section using the audited support package, retaining
all prespecified comparisons, failed or infeasible cases and uncertainty. Until
such evidence exists and has been assessed, this section must remain unfilled.}

\section{Discussion}
\pending{Interpret only verified measured outcomes, including contradictions,
null findings, missing data and feasibility failures. Compare with verified
prior work and avoid converting a surrogate reduction into a transport claim.}
\section{Limitations}
The intended model omits decoherence, many-electron interactions, contacts beyond
ideal leads, fabrication physics and a calibrated voltage-to-potential mapping.
Discrete-band, energy-window, KL-truncation and perturbative errors must be
assessed separately. Conditional fidelity alone can hide negligible throughput.
\pending{Add observed numerical and statistical limitations, including finite
sample uncertainty, conditional inference and any failed convergence checks.}
\section{Conclusion}
\pending{Write only conclusions supported by the completed study. No conclusion
about superiority, novelty, publication readiness or hardware is supplied.}
\bibliographystyle{plain}
\bibliography{verified_references}
\end{document}
